%!TEX TS-program = xelatex
%!TEX encoding = UTF-8 Unicode
%!TEX outputDirectory = output/en
\documentclass[11pt, a4paper]{awesome-cv}
\hypersetup{pageanchor=true}

\geometry{left=1.4cm, top=.8cm, right=1.4cm, bottom=1.8cm, footskip=.5cm}
\fontdir[fonts/]
\colorlet{awesome}{awesome-red}
\setbool{acvSectionColorHighlight}{true}
\renewcommand{\acvHeaderSocialSep}{\quad\textbar\quad}

% Toggle public version (hides sensitive info like phone/email)
\newif\ifpublic
\ifdefined\makePublic \publictrue \fi

% === Shared Info ===
\name{Bettahar}{Samir}
\ifpublic
\else
  \mobile{(+213)~000~00~00~00}
  \email{Bettahar.Samir@outlook.com}
\fi
\github{samir1498}
\linkedin{samir-bettahar}
\homepage{samir-bettahar.dev}

\begin{document}

%========================
% ENGLISH VERSION
%========================
% Fix redefinition issues
\makeatletter
\let\drawphoto\relax
\let\headertextwidth\relax
\let\headerphotowidth\relax
\makeatother

\position{Software Engineer}
% \address{City, Country}

\makecvheader[C]
\makecvfooter{\today}{Samir Bettahar~~~·~~~CV (EN)}{\thepage}

% --- Summary ---
\cvsection{Summary}
\begin{cvparagraph}

Software engineer with strong expertise in JavaScript and TypeScript, experienced in developing modern frontend applications using React, Astro, and SvelteKit. Skilled in building responsive UIs, integrating REST APIs, and ensuring clean, maintainable code. Quick to learn and adapt to new frameworks such as Vue.js, with additional experience in backend development as a complement.
\end{cvparagraph}

\cvsection{Experience}
\begin{cventries}

  \cventry
    {Full-Stack Developer}
    {Omnivya}
    {Remote}
    {Feb. 2024 -- Aug. 2025}
    {
      \begin{cvitems}
        \item {Developed RESTful APIs with \textbf{Spring Boot}, Spring Security, and Spring Data JPA.}
        \item {Implemented data persistence and schema migrations with PostgreSQL and Flyway.}
        \item {Improved performance with caching using Redis.}
        \item {Refactored and built web interfaces with React, AstroJS, and TanStack Router.}
        \item {Implemented authentication solutions (Zitadel, BetterAuth).}
        \item {Wrote unit, integration, and end-to-end tests using Vitest, TestContainers, and Playwright.}
      \end{cvitems}
    }

\end{cventries}

% --- Education ---
\cvsection{Education}
\begin{cventries}

  \cventry
    {Master's in Software Engineering}
    {University of Saad Dahlab Blida}
    {Blida, Algeria}
    {2020 -- 2022}
    {
      \begin{cvitems}
        \item {Focus on software architecture, development methodologies, and full-stack development}
      \end{cvitems}
    }

  \vspace{1em}

  \cventry
    {Bachelor’s in Information Systems and Software Engineering}
    {University of Saad Dahlab Blida}
    {Blida, Algeria}
    {2016 -- 2020}
    {
      \begin{cvitems}
        \item {Foundation in computer science fundamentals, algorithms, databases, and object-oriented programming.}
      \end{cvitems}
    }

\end{cventries}

%========================
% PROJECTS
%========================
\cvsection{Projects}

\begin{cventries}
    \cventry
    {Bachelor Project – Parkinson’s Disease Detection (Python, Machine Learning)}
    {Saad Dahlab University of Blida}
    {Blida, Algeria}
    {2020}
    {
        \begin{cvitems}
            \item {Developed a medical decision support system in \textbf{Python} for detecting Parkinson’s disease using \textit{spiral} and \textit{wave} handwriting tests.}
            \item {Implemented \textbf{machine learning models} (KNN, SVM, Random Forest) with \textbf{scikit-learn}, achieving up to \textbf{90\% accuracy}.}
            \item {Applied image preprocessing and feature extraction techniques with \textbf{OpenCV}, including LBP, Haralick, and Hu moments.}
            \item {Built a \textbf{Tkinter GUI} enabling doctors to test images and visualize confusion matrices and ROC curves.}
        \end{cvitems}
    }
    
    \cventry
    {Master Thesis – Software Product Line for e-Banking Applications (Java EE, Ontologies)}
    {Saad Dahlab University of Blida}
    {Blida, Algeria}
    {2022}
    {
        \begin{cvitems}
            \item {Designed and developed a \textbf{software product line} for automatically generating e-Banking applications based on \textbf{ontologies} and feature models.}
            \item {Implemented a desktop configurator in \textbf{Java Swing} for selecting features and generating applications.}
            \item {Developed full web applications (Client and Admin portals) with \textbf{Java EE} (Servlets, JSP, JDBC) and deployed on Tomcat/Glassfish.}
        \end{cvitems}
    }
\end{cventries}
\newpage
% --- Skills ---
\cvsection{Skills}
\begin{cvskills}

  \cvskill {Languages} {TypeScript, JavaScript, Java, Go}
  \cvskill {Backend} {NestJS, Spring Boot, Go (API for Three.js), REST APIs}
  \cvskill {Frontend} {React.js, AstroJS, SvelteKit, Three.js}
  \cvskill {Database} {PostgreSQL, MongoDB, Spring Data JPA, TypeORM, Flyway}
  \cvskill {Infra \& DevOps} {Docker, CI/CD, Git/GitHub, GraalVM}
  \cvskill {Authentication} {Zitadel, Clerk, BetterAuth}
  \cvskill {Testing} {Vitest, TestContainers, Playwright}
  \cvskill {Observability} {Prometheus, OpenTelemetry}

\end{cvskills}

\end{document}
